\documentclass[12pt]{article}
\usepackage[utf8]{inputenc}
\usepackage[T1]{fontenc}
\usepackage{hyperref}
\title{Wantoat Language Resources for K-12 Curriculum Development}
\author{Compiled by Claude Sonnet 5.5 (claude-sonnet-5-5) for Andr\'as Kornai}
\date{}
\begin{document}
\maketitle

\section{Introduction}
Wantoat (ISO 639-3 wnc) is a Finisterre language of Morobe Province, Papua New Guinea. The figure of 8,200 speakers is cited from 1978. Dialects include Wapu (Leron), Central Wantoat, Bam and Yagawak (Kandomin). Wikipedia notes a UNESCO rating of critically endangered. No school use was found.

\section{Primary Educational Resources}
None found.

\section{Linguistic and Reference Materials}
\begin{itemize}
\item Wikipedia, Wantoat language: overview, dialects and consonant inventory. \url{https://en.wikipedia.org/wiki/Wantoat_language}
\item Omniglot, Wantoat (Taap): alphabet and pronunciation notes credited to Wolfram Siegel. \url{https://omniglot.com/writing/wantoat.htm}
\item Glottolog reference: Davis, Donald R., compiler, 2008 [1958--1974], a Wantoat--English--Tok Pisin dictionary. \url{https://glottolog.org/resource/reference/id/301481}
\item Endangered Languages Project entry. \url{https://www.endangeredlanguages.com/lang/3332?hl=fr}
\end{itemize}
No grammar was found in the search results.

\section{Religious and Community Resources}
Sources disagree on the extent of Scripture. Omniglot says parts of the Bible were published in 2001, while Joshua Project says the New Testament has been translated. \url{https://joshuaproject.net/index.php/people_groups/15858}

\section{Digital and Audio Resources}
Nothing specific was found beyond the entries above.

\section{Curriculum Development Context}
The existing dictionary, an orthography and Scripture portions give a base for vernacular literacy materials. No primers or school programmes were found.

\end{document}